\chapterdivider{2}{Continuation fundamentals}
  {How do we follow a connected family of solutions instead of solving isolated points?}
  {Featured application: the scalar cubic branch and its two folds}

\begin{frame}{Continuation answers a different question}
\centering
\begin{tikzpicture}[>=Latex,node distance=8mm,
  card/.style={rounded corners,draw=JaxBlue,thick,fill=SoftBlue,
    minimum width=40mm,minimum height=28mm,align=center,font=\small}]
  \node[card] (sim) {\textbf{Simulation}\\[1mm]
    start from one state\\follow time $t$};
  \node[card,right=of sim] (root) {\textbf{One root}\\[1mm]
    fix $p$\\solve $F(\state,p)=0$};
  \node[card,right=of root,draw=JaxTeal,fill=JaxTeal!9] (cont)
    {\textbf{Continuation}\\[1mm]
    follow connected pairs\\$(\state,p)$ with $F=0$};
  \draw[->,thick,JaxOrange] (sim)--node[above,font=\scriptsize]{different question}(root);
  \draw[->,thick,JaxOrange] (root)--node[above,font=\scriptsize]{reuse locality}(cont);
\end{tikzpicture}

\vfill
Continuation moves through solution space; it does not integrate physical time.
\end{frame}

\begin{frame}{A branch is a geometric curve}
\centering
\begin{tikzpicture}[x=1.45cm,y=1.05cm,>=Latex]
  \draw[->] (-3.3,0)--(3.3,0) node[right] {$p$};
  \draw[->] (0,-2.7)--(0,2.7) node[above] {$u$};
  \draw[JaxBlue,very thick,smooth,domain=-2.1:2.1,samples=100]
      plot ({\x^3/3-\x},{\x});
  \foreach \x in {-1.8,-1.3,-.8,-.3,.2,.7,1.2,1.7}
    \fill[JaxOrange] ({\x^3/3-\x},{\x}) circle (1.6pt);
  \node[align=left,anchor=west] at (1.4,2.0)
    {ordered solution pairs\\sample one connected curve};
\end{tikzpicture}
\vspace{-2mm}
\[
  F(u,p)=p+u-u^3/3=0.
\]
The branch need not be a single-valued function $u(p)$.
\end{frame}

\begin{frame}{Nearby solutions make continuation efficient}
\centering
\begin{tikzpicture}[>=Latex,node distance=15mm,
  flowbox/.style={rounded corners,draw=JaxBlue,thick,fill=SoftBlue,
    minimum width=38mm,minimum height=19mm,align=center}]
  \node[flowbox] (known) {accepted solution\\$(\state_k,p_k)$};
  \node[flowbox,right=of known,draw=JaxOrange,fill=JaxOrange!10] (guess)
    {predict nearby\\use local direction};
  \node[flowbox,right=of guess,draw=JaxGreen,fill=JaxGreen!10] (exact)
    {correct to $F=0$\\accept or retry};
  \draw[->,very thick,JaxOrange] (known)--(guess);
  \draw[->,very thick,JaxGreen!70!black] (guess)--(exact);
\end{tikzpicture}

\vfill
On a smooth branch, yesterday's solution is already close to today's root.
\end{frame}

\begin{frame}{The Jacobian measures local sensitivity}
Differentiate the equilibrium equation along a regular branch:
\[
  F(\state(p),p)=0
  \quad\Longrightarrow\quad
  F_{\state}\frac{d\state}{dp}+F_p=0,
  \qquad
  \frac{d\state}{dp}=-F_{\state}^{-1}F_p.
\]
\begin{block}{Read the formula geometrically}
When $F_{\state}$ is well conditioned, a small parameter change produces a
controlled state change. As $F_{\state}$ approaches singularity, the
parameter coordinate becomes increasingly sensitive.
\end{block}
\end{frame}

\begin{frame}{At a fold, the parameter turns around}
\centering
\begin{tikzpicture}[x=1.45cm,y=1.05cm,>=Latex,
  note/.style={fill=white,inner sep=2pt,align=center,font=\small}]
  \draw[->] (-3.2,0)--(3.2,0) node[right] {$p$};
  \draw[->] (0,-2.7)--(0,2.7) node[above] {$u$};
  \draw[JaxBlue,very thick,smooth,domain=-2.1:2.1,samples=100]
      plot ({\x^3/3-\x},{\x});
  \fill[JaxRed] ({2/3},{-1}) circle (3pt) node[right=2mm] {fold};
  \fill[JaxRed] ({-2/3},{1}) circle (3pt) node[left=2mm] {fold};
  \draw[JaxOrange,very thick,->]
    (-1.3,-1.75)..controls(-.1,-1.4) and (.8,-1.2)..({2/3},{-1});
  \draw[JaxOrange,very thick,->]
    ({2/3},{-1})..controls(.5,-.6) and (.1,-.2)..(-.2,.3);
  \node[note] (turn) at (1.55,-1.75) {$p$ reverses direction};
  \draw[->,JaxRed] (turn.west)--(.76,-1.18);
  \node[note] (smooth) at (-1.45,.2) {the curve stays smooth};
  \draw[->,JaxBlue] (smooth.south east)--(-.38,-.15);
\end{tikzpicture}
\end{frame}

\begin{frame}{What to remember}
\begin{columns}[T,onlytextwidth]
\column{.48\textwidth}
\begin{block}{Curve viewpoint}
Continuation follows connected roots in the augmented space $(\state,p)$ and
uses nearby accepted points to predict the next solution.
\end{block}

\column{.48\textwidth}
\begin{block}{Fold viewpoint}
At a fold, $F_{\state}$ is singular and $p$ stops being a valid local
coordinate, even though the solution curve can remain smooth.
\end{block}
\end{columns}

\vfill
\centering
\large\textcolor{JaxTeal}{A fold breaks the parameter coordinate, not the branch.}
\end{frame}
